\subsection{矩阵系数}

定义 5. --- 设 G 是拓扑群，\(\varrho\) 是 G 在希尔伯特空间 E 中的酉表示。设 \(x\) 和 \(y\) 是 E 中的元素。由 \(g \mapsto \langle x | \varrho(g)y \rangle\) 给出的从 G 到 K 的函数称为 \(\varrho\) 的矩阵系数，也称代表函数。若 \(x = y\)，则称其为对角矩阵系数。若 \(\varrho\) 是有限维的，则称其为有限维矩阵系数。

\(\varrho\) 的矩阵系数是 G 上的连续有界函数。我们用 \(\Upsilon(G)\)（分别用 \(\Theta(G)\)）表示 G 的复酉表示（分别为复有限维表示）的矩阵系数集合。

命题 5. — 设 G 是拓扑群。集合 \(\Theta(\mathrm{G})\) 和 \(\Upsilon(\mathrm{G})\) 是 \(\mathcal{C}_{\mathrm{b}}(\mathrm{G})\) 的含单位对合子代数。

常值函数 1 是 G 在 C 上的平凡表示的矩阵系数。设 \(\varrho\) 是 G 的酉表示，且 \(x\) 和 \(y\) 是 \(\varrho\) 的表示空间中的向量。对每个 \(\lambda \in \mathbf{C}\) 及每个 \(g \in \mathrm{G}\)，有 \(\lambda \langle x \mid \varrho(g)y \rangle = \langle x \mid \varrho(g)(\lambda y) \rangle\)。此外，有

\[
\langle x \mid \varrho (g) y \rangle = \langle \varrho (g ^ {- 1}) x \mid y \rangle = \overline {{\langle y \mid \varrho (g ^ {- 1}) x \rangle}}
\]

对所有 \(g \in G\) 成立。因此，集合 \(\Theta(G)\) 和 \(\Upsilon(G)\) 在标量乘法和共轭下稳定。

设 \(\varrho_{1}\) 和 \(\varrho_{2}\) 是 G 的酉表示；设 \((x_{1}, y_{1})\) 是 \(\varrho_{1}\) 的表示空间中的向量，\((x_{2}, y_{2})\) 是 \(\varrho_{2}\) 的表示空间中的向量。于是对每个 \(g \in G\)，有

\[
\begin{array}{c} \langle x _ {1}   |   \varrho_ {1} (g) y _ {1} \rangle + \langle x _ {2}   |   \varrho_ {2} (g) y _ {2} \rangle = \langle (x _ {1}, x _ {2})   |   (\varrho_ {1} \oplus \varrho_ {2}) (g) (y _ {1}, y _ {2}) \rangle , \\ \langle x _ {1}   |   \varrho_ {1} (g) y _ {1} \rangle \langle x _ {2}   |   \varrho_ {2} (g) y _ {2} \rangle = \langle x _ {1} \otimes x _ {2}   |   (\varrho_ {1} \otimes \varrho_ {2}) (g) (y _ {1} \otimes y _ {2}) \rangle , \end{array}
\]

这证明了 \(\Theta(G)\) 和 \(\Upsilon(G)\) 在加法和乘法下稳定。命题得证。

